\documentclass[10pt,a4paper]{amsart}
\usepackage[margin=19mm]{geometry}
\usepackage{amsmath,amssymb,mathtools}
\usepackage[scr=rsfs,cal=euler]{mathalfa}
\usepackage{xfrac}
\usepackage[unicode,hidelinks,bookmarks=false]{hyperref}
\newcommand{\ie}{i.e.\ }
\newcommand{\cf}{cf.\ }
\newcommand{\resp}{respectively\ }
\newcommand{\Subsection}{\subsection}
\providecommand{\nobreakdash}{\nobreak\hbox{-}}
\setlength{\emergencystretch}{2em}
\multlinegap0pt
\usepackage{bm}
\usepackage{bbm}
\usepackage{dsfont}
\usepackage{xspace}
\usepackage{cancel}
\usepackage{mathtools}
\usepackage{cleveref}
\usepackage{amsthm}
\makeatletter
\@ifundefined{lemma}{\newtheorem{lemma}{Lemma}}{}
\@ifundefined{theorem}{\newtheorem{theorem}{Theorem}}{}
\makeatother
\title[Borderline regularity in singular free boundary problems]{Borderline regularity in singular free boundary problems}
\author{Dami\~ao J. Ara\'ujo, Aelson Sobral, Eduardo V. Teixeira, Jos\'e Miguel Urbano}
\date{Source: arXiv:2508.14736v1 (CC BY 4.0)}
\begin{document}
\maketitle

\noindent\emph{Source and licence.} Borderline regularity in singular free boundary problems. Version arXiv:2508.14736v1 (CC BY 4.0), distributed under the Creative Commons Attribution 4.0 International licence, as recorded in the arXiv licence metadata for this exact version and shown on the abstract page https://arxiv.org/abs/2508.14736v1. Full rights notice: licenses/arxiv-2508.14736-CC-BY-4.0.txt.\par\smallskip

\noindent\textit{Editorial scope.} Counted unit: the Theorem whose source label is \texttt{C1 regularity at free boundary} in the article (source0.tex lines 596--602, source label \texttt{C1 regularity at free boundary}), delivered together with the article's own complete original proof of it (source0.tex lines 604--720, one \texttt{proof} environment of 117 lines). No mathematics is written, repaired or re-derived: statement and proof are the article's own text line for line. Required context retained uncounted: main functional (lines 152--154); flat lemma (lines 517--526). Every definition, display and label of those lines is retained unchanged. Omitted from this package and not redistributed: 1--151, 155--516, 527--595, 603--603, 721--860. Nothing inside a retained excerpt is omitted or altered: every retained body is delivered exactly as the article's lines are written, so no omission receipt of any kind accompanies this package. Citations: the retained excerpts carry no citation command at all, so no bibliography entry of the article is transcribed and no reference is redistributed. Reference adaptation: none was needed and none was made; every reference inside the delivered unit resolves inside it, so no unresolved reference or citation is produced. Printed slips kept literal: no printed word, symbol, hypothesis or number of the counted statement or of its proof was altered. Layout and class: the article's own class is \texttt{amsart} with packages including amsmath, amssymb, amsthm, mathtools, mathrsfs, hyperref, cleveref, xcolor, tikz, verbatim, setspace, ulem; the wrapper is the lane's pinned \texttt{amsart} editorial wrapper at 10pt on A4, which restates the theorem-like environments the retained text needs and adds the article's own macro definitions that the retained text needs (none), with extra wrapper packages (bm, bbm, dsfont, xspace, cancel, mathtools, cleveref). The delivered numbering is the unit's own. Assets: no retained span contains a figure environment, a graphics-inclusion command, a PDF-page inclusion, a table or a TikZ picture; no image file is copied into this package, the package declares no asset dependency and no image file is redistributed with it. Licence: version arXiv:2508.14736v1 of this article is distributed under the Creative Commons Attribution 4.0 International licence, as recorded in the arXiv licence metadata for this exact version and shown on the abstract page https://arxiv.org/abs/2508.14736v1. A full rights notice is delivered at licenses/arxiv-2508.14736-CC-BY-4.0.txt. Characters: every retained excerpt is pure ASCII, so the delivered engine, XeLaTeX with its default Latin Modern text font, reports no missing character.
\begin{equation}\label{main functional}
\mathcal{J}(u,\Omega) \coloneqq \int_{\Omega} \frac{1}{2}|Du|^2 + \sigma(u)\,dx,
\end{equation}

\begin{lemma}\label{flat lemma}
Given $\epsilon > 0$, there is $\delta>0$, depending only on $\epsilon$, such that for any $u$ nonnegative local minimizer of \eqref{main functional} in $B_1$, with 
$$
   u(0) = 0, \quad 0 \leq u \leq 1, \text{ and }\, \|\sigma\|_{L^\infty(\mathbb{R})} \leq \delta,
$$
there holds
$$
    \sup_{B_{1/2}} u \leq \epsilon.
$$
\end{lemma}

\begin{theorem}\label{C1 regularity at free boundary}
Let $u$ be a nonnegative local minimizer of \eqref{main functional} in $B_1$, with $\sigma$ being a modulus of continuity. Then, there exists a modulus of continuity $\omega \colon [0,1] \to [0,+\infty)$, depending on universal parameters and $\sigma$, such that, for any $x_0 \in \partial \{u>0\}\cap B_{1/2}$, 
$$
    |u(x)| \leq C\left(\|u\|_{L^\infty(B_{3/4})} + \sqrt{\|\sigma\|_{L^\infty(\mathbb{R})}} \, \right)|x-x_0|\,\omega(|x-x_0|), \quad x \in B_{1/4}(x_0),
$$
for some universal constant $C>0$.
\end{theorem}

\begin{proof}
Let $x_0 \in \partial \{u>0\} \cap B_{1/2}$ be an arbitrary point. Let $\epsilon = 2^{-2}$ in \Cref{flat lemma}, and consider the corresponding $\delta >0$. Up to replacing $u$ by
$$
    v(x) \coloneqq \frac{u(x_0 + 2^{-1}x)}{\|u\|_{L^\infty(B_{3/4})} + \sqrt{\delta^{-1}\|\sigma\|_{L^\infty(\mathbb{R})}}}, \quad x \in B_1,
$$
we may assume $u$ satisfies the assumption of \Cref{flat lemma} for this specific $\epsilon$.

We will prove that there is a sequence $(\mu_k)_{k \in \mathbb{N}} \subset (0,1)$ such that
$$
    \sup_{B_{2^{-k}}}u \leq 2^{-k}\left(\prod_{i=1}^{k}\mu_i\right), \quad \text{for every }\, k \in \mathbb{N}.
$$
This sequence will be built via an induction process. The first step, as we already are in a normalized setting, is to apply \Cref{flat lemma} to obtain
\begin{equation}\label{normalization step 1}
    \sup_{B_{2^{-1}}} u \leq  2^{-2} \leq 2^{-1} \mu_1,
\end{equation}
for some $2^{-1} \leq \mu_1 < 1$ to be chosen. To fix it, define the function
$$
    v_1(x) \coloneqq \frac{u\left(2^{-1}x\right)}{\mu_1 2^{-1}}, \quad x \in B_1,
$$
which is normalized due to \eqref{normalization step 1}. Moreover, $v_1$ is a local minimizer to
$$
    \mathcal{J}_1(w, B_1) \coloneqq \int_{B_1}\frac{1}{2}|Dw|^2 + \sigma_1(w)\, dx, \quad \text{with} 
    \quad \sigma_1(t) \coloneqq \mu_1^{-2}\sigma\left(\mu_1 2^{-1} t\right).
$$
At this point, we desire to apply \Cref{flat lemma} to $v_1$. As $v_1(0) = 0$, the only missing assumption is the smallness condition for $\sigma_1$. As $\sigma_1$ is a modulus of continuity, the smallness condition translates to ensuring $\sigma_1(1) \leq \delta$. This, in turn, is equivalent to
$$
    \mu_1^{-2}\sigma\left(\mu_1 2^{-1}\right) \leq \delta.
$$
To make this selection in a proper way, we build the following algorithm: if 
$$
   2^2\sigma\left(2^{-2}\right) \leq \delta,
$$
we pick $\mu_1 = 2^{-1}$. If this does not happen, then we pick $1 > \mu_1 > 2^{-1}$ so that
$$
    \mu_1^{-2}\sigma(\mu_1 2^{-1}) = \delta.
$$
Indeed, define
$$
    f(t) \coloneqq t^{-2} \sigma(2^{-1}t), \quad \text{for }\, t \in [2^{-1},1].
$$
This is a continuous function, such that 
$$f(1) = \sigma(2^{-1}) < \sigma(1) \leq \delta < f(2^{-1}),$$
and so, by the Intermediate Value Theorem, there should be $1 > \mu_1 > 2^{-1}$ such that 
$$
    f(\mu_1) = \delta.
$$
Now, we can apply \Cref{flat lemma} to $v_1$ to obtain
$$
    \sup_{B_{2^{-1}}} v_1 \leq 2^{-2} \ \Longrightarrow \ \sup_{B_{2^{-2}}} u \leq 2^{-2} \mu_1 \mu_2,
$$
for $\mu_1 \leq \mu_2 < 1$ to be chosen. Next, define
$$
    v_2(x) \coloneqq \frac{v_1\left(2^{-1}x\right)}{\mu_2 2^{-1}}, \quad x \in B_1.
$$
Again, it follows that $v_2$ is normalized, $v_2(0) = 0$ and it is a local minimizer of
$$
    \mathcal{J}_2(w,B_1) \coloneqq \int_{B_1} \frac12 |Dw|^2 + \sigma_2(w)dx, \quad \text{with} \quad \sigma_2(t) \coloneqq \mu_2^{-2}\sigma_1\left(\mu_2 2^{-1} t\right).
$$
We then follow the same routine: if
$$
    \mu_1^{-2}\sigma_1\left(\mu_1 2^{-1} \right) \leq \delta,
$$
then we pick $\mu_2 = \mu_1$. If not, then we pick $1 > \mu_2 > \mu_1$ so that
$$
   \mu_2^{-2}\sigma_1\left(\mu_2 2^{-1}\right) = \delta. 
$$
We can repeat this process inductively: we assume we have built up to $k$, that is, we have already found the parameter $\mu_k$ and the functions $\sigma_k$ and $v_k$ are already defined. For a parameter $1 > \mu_{k+1} \geq \mu_k$, we define the function
$$
    v_{k+1}(x) \coloneqq \frac{v_{k}(2^{-1}x)}{\mu_{k+1}2^{-1}}, \quad x \in B_1,
$$
which is normalized, $v_{k+1}(0) = 0$ and minimizes a functional with $\sigma_{k+1}(t) \coloneqq \mu_{k+1}^{-2}\sigma_k(\mu_{k+1}2^{-1}t)$. The choice of $\mu_{k+1}$ is as follows: if
$$
    \mu_{k}^{-2}\sigma_k(\mu_{k}2^{-1}) \leq \delta,
$$
then we pick $\mu_{k+1} = \mu_k$. If not, then we pick $\mu_{k+1}$ such that
$$
    \mu_{k+1}^{-2}\sigma_k\left(\mu_{k+1} 2^{-1}\right) = \delta.
$$
This concludes the construction of the sequence of parameters $\mu_k$, and so we have proven that
$$
    \sup_{B_{2^{-k}}}u \leq 2^{-k}\left(\prod_{i=1}^{k}\mu_i\right), \quad \text{for every }\, k \in \mathbb{N}.
$$
Define
$$
    a_k \coloneqq \prod_{i=1}^{k}\mu_i.
$$
By construction, $(a_k)_{k \in \mathbb{N}} \subset (0,1)$ and is decreasing. As a consequence, $a_k \to a_{\ast}$ for some $a_{\ast} \in [0,1)$. We prove that $a_{\ast} = 0$. Indeed, if after some point, the sequence stabilizes, meaning that
$$
    \mu_k = \mu_{k_0}, \quad \text{for all }\, k \geq k_0,
$$
then it readily follows that $a_{\ast} = 0$. Otherwise, we would have, for some subsequence $(k_j)_{j \in \mathbb{N}}$, that
$$
    a_{k_j}^{-2}\sigma(a_{k_j}2^{-k_j}) = \delta.
$$
Since $a_k \to a_{\ast}$, it also follows that $a_{k_j} \to a_{\ast}$, and so, passing to the limit in the equation above, gives
$$
    \sigma(a_{\ast} \cdot 0) = \delta a_{\ast}^2.
$$
But since $\sigma(0) = 0$, it follows that $a_{\ast} = 0$. 

From here, we can build the modulus of continuity as follows: pick $x \in B_{2^{-2}}$ and let $\rho = |x|$. There exists $k \in \mathbb{N}$ such that $2^{-(k+1)} < \rho \leq 2^{-k}$. Therefore, 
$$
    |u(x)| \leq \sup_{B_{\rho}} u \leq \sup_{B_{2^{-k}}}u \leq 2^{-k} a_k \leq 2\rho a_k,
$$
since $2^{-k} \leq 2\rho$. Taking into account that
$$
    k \leq \frac{\ln\left(\rho^{-1}\right)}{\ln(2)} < k+1,
$$
we may then define
$$
     \omega(\rho) \coloneqq 2 \prod_{i=1}^{\left\lfloor{\frac{\ln(\rho^{-1})}{\ln(2)}}\right\rfloor} \mu_i.
$$
Observe that $\omega$ is an increasing function and 
$$
    \lim_{\rho \to 0} \omega(\rho) = 2 \lim_{k \to \infty} a_k = 0.
$$
\end{proof}

\end{document}
